\documentclass[11pt,a4paper]{article}
\usepackage[utf8]{inputenc}
\usepackage[T1]{fontenc}
\usepackage{lmodern,amsmath,amssymb,textcomp}
\usepackage[margin=24mm]{geometry}
\usepackage{longtable,booktabs,array,enumitem}
\usepackage[hidelinks]{hyperref}
\setlength{\parindent}{0pt}
\setlength{\parskip}{6pt}
\emergencystretch=3em
\hypersetup{pdfauthor={Rachel Teo, Wenyi Wang, Hamdi Abderrahmene, Alejandro Zarzuelo Urdiales}}
\begin{document}
{\LARGE\bfseries Arbitrarily long prime-exponent gaps in complete-interval subset products\par}\medskip
{\small Rachel Teo\ensuremath{^1}, Wenyi Wang\ensuremath{^1}, Hamdi Abderrahmene\ensuremath{^1}, Alejandro Zarzuelo Urdiales\ensuremath{^2}\par}\medskip
{\small\textbf{Affiliations:} \href{https://sakana.ai/}{\ensuremath{^1} Sakana AI}; \href{https://rsihouse.ai/openmath}{\ensuremath{^2} Open Math}.\par}

{\small\href{https://github.com/alejandrozu/openmath-2026-judging}{Code and formal proofs}\par}

\subsection{Abstract}
The distinct-subset-product interpolation question associated with OEIS A060957 is refuted. More strongly, for every positive gap length and every prescribed lower bound on a prime, the construction realizes an exact fibre consisting of two blocks of four attainable exponents separated by the prescribed missing block.

\subsection{The original interpolation problem}
Let P(n) be the set of products of subsets of \{1,...,n\}. A factor may be chosen once or omitted: this is not the unrestricted multiplicative monoid generated by that interval. The question asks whether, when m and p\textasciicircum{}a m are attainable for a prime p\ensuremath{\leq}n, all the intermediate p-exponent levels are attainable. Its complete-interval condition matters. A hole in a hand-picked collection of factors would not answer the question, because the omitted interval factors could fill that hole.

The submitted result handles all factors of the interval simultaneously. Its basic endpoint supplies m, p and n with m and p\textasciicircum{}5m in P(n), but p\textasciicircum{}4m outside P(n). The stronger endpoint specifies the entire exponent fibre, rather than just proving that one level is missing. For every g\ensuremath{\geq}1 and lower bound B, there are n,p,d,h with p prime, B<p\ensuremath{\leq}n, d>0 and p not dividing d, for which p\textasciicircum{}e d is attainable precisely at e=h,h+1,h+2,h+3,h+g+4,h+g+5,h+g+6,h+g+7. The g omitted levels are h+4 through h+g+3.

\subsection{Encoding and the exposed-face construction}
The proof represents integer products by their prime-exponent vectors, separating the distinguished p coordinate from all other coordinates. A graph-based atom construction controls the possible exponent increments. In the generalized construction, a circulant system decomposes the complete graph on an odd number of vertices into oriented spanning two-factors. Every vertex sees every other vertex once through the shifts and their inverses. This identity is proved from the circulant model; it is not left as an existence hypothesis in the final theorem.

Integral feasibility conditions in the height-one slice classify the allowable atoms. Summing inequalities by colour gives a degree-two budget; summing by vertex gives the complementary bound. These estimates constrain integral solutions to the two desired exponent blocks. The argument is algebraic and combinatorial rather than a search through arbitrary integer factors.

To transport that local construction to the complete interval, the proof chooses prime coordinates and an isolating rational linear functional. The positive-weight factors form a forced baseline W. Any subset product with the same p-free exponent vector must contain these positive-weight factors and can differ only through zero-weight atoms in the prescribed subspace. Consequently unrelated factors from \{1,...,n\} cannot silently create extra attainable levels. The complete-face theorem proves both directions of this classification. Finally, dividing the baseline by its p-primary factor produces d not divisible by p and an absolute exponent offset h.

\subsection{Formal theorem and scope}
The passage from a model fibre to products of actual distinct interval factors is proved by the complete-face, product-support and normalization modules. The Main and Generic.Final endpoints establish the explicit counterexample, failure of interpolation, exact arbitrary-gap eight-level fibre and consecutive-above corollary.

The ordinary-proof expansion now gives the complete atom classification, the choice of prime parameters and separating linear functional, and the bridge preventing the remaining factors of the complete interval from filling the gap. The formal proof supplies a precise supplementary specification alongside that human-readable argument. No minimum counterexample size is asserted, and the theorem does not provide one fixed prime that works for every gap length.

\subsection{Prior work and significance}
Prime-coordinate encoding, exposed faces and semigroup holes are established methods. Here they yield a counterexample to complete-interval interpolation and gaps of arbitrary length with an exact eight-level fibre. Historical priority for this precise combination remains unresolved after the bounded comparison with the cited sources.

\subsection{Formal verification}
Lean 4.33.1 checked the 45-source development and all four selected endpoint declarations. Their axiom dependencies are confined to propext, Classical.choice and Quot.sound; none uses an admitted proof, native evaluation or an unknown axiom. Exact source pins, proof files and audit records are supplied in the companion repository.

\begin{samepage}
\subsection{SubsetProductGap.Generic.arbitrary\_gap\_eight\_levels\_above}
\begin{quote}\small\ttfamily theorem arbitrary\_gap\_eight\_levels\_above (g : \ensuremath{\mathbb{N}}) (hg : 1 \ensuremath{\leq} g) (lower : \ensuremath{\mathbb{N}}) :\newline     \ensuremath{\exists} n p d h : \ensuremath{\mathbb{N}}, p.Prime \ensuremath{\wedge} lower < p \ensuremath{\wedge} p \ensuremath{\leq} n \ensuremath{\wedge} 0 < d \ensuremath{\wedge} \ensuremath{\neg} p \ensuremath{\mid} d \ensuremath{\wedge}\newline       \ensuremath{\forall} e : \ensuremath{\mathbb{N}}, p\textasciicircum{}e * d \ensuremath{\in} productsOfSubsets n \ensuremath{\leftrightarrow}\newline         e = h \ensuremath{\vee} e = h + 1 \ensuremath{\vee} e = h + 2 \ensuremath{\vee} e = h + 3 \ensuremath{\vee}\newline           e = h + g + 4 \ensuremath{\vee} e = h + g + 5 \ensuremath{\vee} e = h + g + 6 \ensuremath{\vee} e = h + g + 7\end{quote}
{\small Source: proofs/a060957/OpHack/SubsetProducts/Generic/Final.lean:90.\par}

{\small Source SHA-256: a2e26af97d551980ecca1eb62f5db1362f75ada87f28795f99fe526651007d53\par}

\end{samepage}
\begin{samepage}
\subsection{SubsetProductGap.Generic.arbitrary\_gap\_consecutive\_above}
\begin{quote}\small\ttfamily theorem arbitrary\_gap\_consecutive\_above (g : \ensuremath{\mathbb{N}}) (hg : 1 \ensuremath{\leq} g) (lower : \ensuremath{\mathbb{N}}) :\newline     \ensuremath{\exists} n p M : \ensuremath{\mathbb{N}}, p.Prime \ensuremath{\wedge} lower < p \ensuremath{\wedge} p \ensuremath{\leq} n \ensuremath{\wedge} 0 < M \ensuremath{\wedge}\newline       M \ensuremath{\in} productsOfSubsets n \ensuremath{\wedge} p\textasciicircum{}(g + 1) * M \ensuremath{\in} productsOfSubsets n \ensuremath{\wedge}\newline         \ensuremath{\forall} j : \ensuremath{\mathbb{N}}, 1 \ensuremath{\leq} j \ensuremath{\to} j \ensuremath{\leq} g \ensuremath{\to} p\textasciicircum{}j * M \ensuremath{\notin} productsOfSubsets n\end{quote}
{\small Source: proofs/a060957/OpHack/SubsetProducts/Generic/Final.lean:113.\par}

{\small Source SHA-256: a2e26af97d551980ecca1eb62f5db1362f75ada87f28795f99fe526651007d53\par}

{\small\textbf{Sources and attribution:} \href{https://oeis.org/A060957}{1. oeis.org / A060957}; \href{https://arxiv.org/abs/2501.02695}{2. arXiv source: 2501.02695}; \href{https://github.com/alejandrozu/openmath-2026-judging/blob/d0ed487f487fa7ec814ff705ab55249983fe8707/OpenMath-Judging/review/entries/E02-a060957.txt}{3. alejandrozu/openmath-2026-judging / E02-a 060957.txt}; \href{https://github.com/alejandrozu/openmath-2026-judging}{Proof source repository}.\par}

\end{samepage}
{\small \textbf{AI assistance.} Codex assisted literature checking, editorial preparation and analysis of the supplied formal artifacts. The human authors are responsible for checking the final mathematical claims, references and attribution.\par}

\end{document}
